\begin{tabular}{|l|c|c|c|c|} \hline
 \bfseries Participant & $\bm{F_{CE}^{max}}$ [N] & $\bm{L_{CE}^{opt}}$ [cm] & $\bm{L_{SEE}^{0}}$ [cm] & \bfseries RMSD [\%] \\ \hline
1 & 4400 & 8.6 & 17 & 7.3 \\ \hline
2 & 5100 & 9.9 & 15 & 6.9 \\ \hline
3 & 5600 & 8.7 & 17 & 5.8 \\ \hline
4 & 6800 & 8.6 & 16 & 4.1 \\ \hline
5 & 9200 & 7.4 & 18 & 4.7 \\ \hline
6 & 8500 & 8.8 & 16 & 2.9 \\ \hline
AVG±SD & 6600 $\pm$ 1700 & 8.7 $\pm$ 0.7 & 17 $\pm$ 1 & 5.3 $\pm$ 1.6 \\ \hline
\end{tabular}\break\hfill\footnotesize{$F_{CE}^{max}$, $L_{CE}^{opt}$ and $L_{SEE}^{0}$ were estimated from the isometric measurements for each participant individually. The root mean square difference (RMSD) was normalised by the individual maximal predicted knee joint moment, which is the product of $F_{CE}^{max}$ and the moment arm of \textit{m. quadriceps femoris.}}